\section{Measuring performance on HLG}
\label{sec:measuring}

\subsection{HLG-RHAE}
\label{sec:rhae}

HLG-RHAE is the per-level Relative Human Action Efficiency metric, structurally
identical to ARC-AGI-3's RHAE \cite{arcagi3}. For each level $l$ in $\{1,\ldots,n\}$ with
$n=34$:

\begin{definition}[Per-level efficiency]\label{def:perlevel}
Let $a_l$ be the agent's move count on its first WIN of level $l$ (or $\infty$ if the
level is never won), and let $h_l$ be the human-baseline move count. The per-level
efficiency score is
\begin{equation}
S_l \;=\; \min\!\left(1.15,\ \frac{h_l}{a_l}\right)^2,
\qquad S_l \coloneqq 0 \ \text{if level $l$ unsolved.}
\label{eq:Sl}
\end{equation}
\end{definition}

\begin{definition}[HLG-RHAE]\label{def:rhae}
With linear weights $w_l = l$ and the set of solved levels $\mathcal{W}$,
\begin{equation}
\mathrm{HLG\text{-}RHAE}
\;=\; \min\!\left(
    \frac{\sum_{l\in\mathcal{W}} w_l}{\sum_{l=1}^{n} w_l},\
    \frac{\sum_{l=1}^{n} w_l \cdot S_l}{\sum_{l=1}^{n} w_l}
  \right).
\label{eq:rhae}
\end{equation}
\end{definition}

The first term in the outer $\min$ is a per-environment cap that ties the maximum
achievable score to the weighted fraction of levels solved. This prevents an agent that
solves only the early easy levels from receiving a high score by being slightly
super-efficient on those levels.

\subsection{Design decisions}
\label{sec:scoring-design}

\paragraph{Linear level weighting.} Level $l$ contributes weight $w_l = l$. Because
optimal solution length grows roughly monotonically in level index, this weight schedule
matches difficulty.

\paragraph{Per-level 1.15$\times$ cap.} If an agent finds a 2-action exploit on a
20-action level, an uncapped ratio would produce $10\times$ credit and overwhelm the
environment average. The cap follows ARC's rationale.

\paragraph{Power-law (squared) efficiency term.} A linear term gives 50\% credit at
$2\times$ human baseline; the squared term gives $25\%$. This sharpens the metric's
discrimination between near-human and substantially worse performance.

\paragraph{Action budget of 5$\times$ optimal.} For a level with optimal length $L_l$,
the agent is terminated after $5 L_l$ moves. This keeps API costs bounded and is
identical in spirit to ARC's $5n$ rule applied to the human-baseline median.

\paragraph{First-success only.} Subsequent attempts on an already-won level do not
modify the score. This mirrors a ``first-contact'' interpretation of generalization.

\subsection{Leaderboards}
\label{sec:leaderboards}

We report two leaderboards.

\paragraph{Official.} The official leaderboard uses the canonical system prompt from
\texttt{system\_prompt.md}, no harness, and the \texttt{none}-effort and \texttt{high}-effort
variants of each model. This produces an apples-to-apples comparison across providers.

\paragraph{Community.} The community leaderboard accepts arbitrary harnesses and
self-reported scores, in the spirit of ARC's community track.

% Auto-generated by scripts/build_paper_assets.py.
% Re-run after extract_leaderboard.py to refresh.
\begin{table}[h]
\centering
\caption{Initial frontier-model HLG-RHAE on the full scoring set (levels 1--33), turn mode, parallel runs. Numbers are computed by \texttt{scripts/extract\_leaderboard.py} against the v0.1 baseline ($h_l = L_l$, optimal length). Cost computed at provider-published per-token rates as of March 2026.}
\label{tab:leaderboard}
\small
\begin{tabular}{@{}llrrrrr@{}}
\toprule
\textbf{Provider} & \textbf{Model} & \textbf{Effort} & \textbf{HLG-RHAE} & \textbf{Won} & \textbf{Tokens} & \textbf{Cost (USD)} \\
\midrule
OpenAI & \texttt{oai-gpt-4.1-2025-04-14} & high & 0.18\% & 1/1 & 23,546 & \$0.12 \\
Anthropic & \texttt{claude-opus-4-6} & high & 0.00\% & 0/1 & 0 & \$0.00 \\
\bottomrule
\end{tabular}
\end{table}

